\documentclass[11pt,a4paper]{article}

% Datos de la actividad
\newcommand{\codigoactividad}{FORM-FB-01}
\newcommand{\nombreactividad}{Fichas bibliográficas: fundamentos de sistemas distribuidos y arquitecturas distribuidas}
\newcommand{\integrantes}{Castro Coello Mario Sebastián}
\newcommand{\repositorio}{https://github.com/reyshadow0/isr701-form-fb-01}
\newcommand{\fechaentrega}{5 de noviembre de 2026}

\usepackage[T1]{fontenc}
\usepackage[utf8]{inputenc}
\usepackage[spanish,es-tabla,es-nodecimaldot]{babel}
\usepackage{lmodern}
\usepackage{microtype}
\usepackage[a4paper,margin=2.5cm]{geometry}
\usepackage{graphicx}
\usepackage{booktabs}
\usepackage{tabularx}
\usepackage{float}
\usepackage{enumitem}
\setlist[description]{style=nextline}
\usepackage{xcolor}
\usepackage{listings}
\usepackage{csquotes}
\usepackage[backend=biber,style=ieee,maxcitenames=6,mincitenames=1]{biblatex}
\addbibresource{referencias.bib}
\usepackage[hidelinks]{hyperref}

\graphicspath{{figuras/}}

% Permite cortar URL largas en la bibliografía (evita desbordes)
\setcounter{biburlnumpenalty}{100}
\setcounter{biburlucpenalty}{100}
\setcounter{biburllcpenalty}{100}

\lstset{
  basicstyle=\ttfamily\small,
  breaklines=true,
  frame=single,
  numbers=left,
  numberstyle=\tiny,
  showstringspaces=false,
  extendedchars=true,
  literate={á}{{\'a}}1 {é}{{\'e}}1 {í}{{\'i}}1 {ó}{{\'o}}1 {ú}{{\'u}}1
           {Á}{{\'A}}1 {É}{{\'E}}1 {Í}{{\'I}}1 {Ó}{{\'O}}1 {Ú}{{\'U}}1
           {ñ}{{\~n}}1 {Ñ}{{\~N}}1 {ü}{{\"u}}1 {¿}{{?`}}1 {¡}{{!`}}1
}
\renewcommand{\lstlistingname}{Código}

\hypersetup{pdftitle={\codigoactividad: \nombreactividad}}

\begin{document}

\begin{titlepage}
  \centering
  {\large\bfseries UNIVERSIDAD TÉCNICA ESTATAL DE QUEVEDO\par}
  \vspace{0.3cm}
  {Facultad de Ciencias de la Computación\par}
  {Carrera de Software\par}
  \vspace{2.5cm}
  {\large Aplicaciones Distribuidas (ISR-701)\par}
  \vspace{1.5cm}
  {\Large\bfseries \codigoactividad\par}
  \vspace{0.4cm}
  {\LARGE \nombreactividad\par}
  \vfill
  \begin{tabular}{@{}rl@{}}
    \textbf{Estudiante:}       &\begin{tabular}[t]{@{}l@{}}\integrantes\end{tabular} \\[2pt]
    \textbf{Docente:}          & Gleiston Cicerón Guerrero Ulloa, PhD \\
    \textbf{Nivel / paralelo:} & Séptimo nivel / B \\
    \textbf{Periodo:}          & SPA 2026--2027 \\
    \textbf{Fecha:}            & \fechaentrega \\
  \end{tabular}
  \par\vspace{1cm}
  {\small Repositorio: \url{\repositorio}\par}
\end{titlepage}

\tableofcontents
\newpage

\section{Fichas bibliográficas}

\subsection{Ficha 1}
\begin{description}
  \item[Referencia (IEEE):] \fullcite{Velepucha2023}
  \item[DOI o ISBN verificado:] \url{https://doi.org/10.1109/ACCESS.2023.3305687} (verificado el 05/11/2026)
  \item[Resumen:]
  \item[Dato concreto:]
  \item[Nota personal:]
\end{description}

\subsection{Ficha 2}
\begin{description}
  \item[Referencia (IEEE):] \fullcite{Ferreira2025}
  \item[DOI o ISBN verificado:] \url{https://doi.org/10.1109/ACCESS.2025.3558923} (verificado el 05/11/2026)
  \item[Resumen:]
  \item[Dato concreto:]
  \item[Nota personal:]
\end{description}

\subsection{Ficha 3}
\begin{description}
  \item[Referencia (IEEE):] \fullcite{Rehman2022}
  \item[DOI o ISBN verificado:] \url{https://doi.org/10.1109/ACCESS.2022.3182211} (verificado el 05/11/2026)
  \item[Resumen:]
  \item[Dato concreto:]
  \item[Nota personal:]
\end{description}

\subsection{Ficha 4}
\begin{description}
  \item[Referencia (IEEE):] \fullcite{Cazzola2022}
  \item[DOI o ISBN verificado:] \url{https://doi.org/10.1007/s00446-022-00429-7} (verificado el 05/11/2026)
  \item[Resumen:]
  \item[Dato concreto:]
  \item[Nota personal:]
\end{description}

\subsection{Ficha 5}
\begin{description}
  \item[Referencia (IEEE):] \fullcite{Blinowski2022}
  \item[DOI o ISBN verificado:] \url{https://doi.org/10.1109/ACCESS.2022.3152803} (verificado el 05/11/2026)
  \item[Resumen:]
  \item[Dato concreto:]
  \item[Nota personal:]
\end{description}

\section{Declaración de uso de inteligencia artificial}

\begin{table}[H]
  \centering
  \caption{Herramientas de inteligencia artificial usadas como apoyo}
  \label{tab:declaracion-ia}
  \begin{tabularx}{\textwidth}{@{}llXX@{}}
    \toprule
    \textbf{Herramienta} & \textbf{Versión} & \textbf{Para qué se usó} & \textbf{Cómo se verificó lo obtenido} \\
    \midrule
     & & & \\
    \bottomrule
  \end{tabularx}
\end{table}

\printbibliography[heading=bibintoc,title={Referencias}]

\end{document}
